\section{The model with constant embodied materials intensity in investment goods}\label{sec:app:constantPhiI}

This is the special case $\Phi^I\!\left(t\right)=\bar\Phi^I$ of
Appendix~\ref{sec:app:exogenousPhiI}, so everything established there applies:
Proposition~\ref{prop:app:exo:zeta} gives $\zeta=0$ wherever $R>\bar\Phi^IG$, the material wedges
vanish, $\Psi=F_R$, and the policy collapses to the two rates
of~\eqref{eq:app:exo:policy}. What constancy adds is a reduction in the dimension of the problem, and
that is the only thing recorded here.

\subsection{The embodied stock ceases to be a state}\label{sec:app:const:state}

\begin{lemma}[$M^K$ is redundant]\label{lem:app:const:state}
Let $Z\equiv M^K-\bar\Phi^IK$. Then $\dot Z=-\delta Z$ along any feasible path,
so $Z\left(t\right)=Z_0e^{-\delta t}$. If the initial data satisfy $M^K_0=\bar\Phi^IK_0$, then
$M^K\left(t\right)=\bar\Phi^IK\left(t\right)$ and $\Phi^K\left(t\right)=\bar\Phi^I$ at every date, and
the ledger~\eqref{eq:app:exo:ledger} reads
\begin{align}
W &= R-\bar\Phi^II+\mathcal N .
\label{eq:app:const:ledger}
\end{align}
\end{lemma}

\begin{proof}
Differentiating $Z$ and using $\dot M^K=\bar\Phi^I\left(I+\delta K\right)-\delta M^K$
from~\eqref{eq:sp:setup:Phi-motion} together with $\dot K=I$,
\begin{align*}
\dot Z &= \bar\Phi^II+\bar\Phi^I\delta K-\delta M^K-\bar\Phi^II \;=\; -\delta\left(M^K-\bar\Phi^IK\right).
\end{align*}
With $Z_0=0$ the stock is $M^K=\bar\Phi^IK$ throughout, so $\dot M^K=\bar\Phi^II$ and the ledger
follows.
\end{proof}

The economics of~\eqref{eq:app:const:ledger} is that replacement investment embodies exactly what
depreciation releases, so only \emph{net} capital formation holds material back from the waste flow.
The planner therefore carries four states $\left(K,S,X,P\right)$ rather than five, and two static
prices $\left(\lambda,p^W\right)$ rather than three. Nothing is lost under inconsistent initial data
either: $Z$ is an exogenous exponential, so~\eqref{eq:app:const:ledger} picks up the known forcing term
$+\delta Z_0e^{-\delta t}$ and the problem stays four-dimensional.

\subsection{The capital margin}\label{sec:app:const:capital}

Eliminating $M^K$ changes what the costate on capital prices, and this is the one place where the
conditions have to be restated rather than read off. With $\bar\Phi^II$ replacing
$\Phi^IG-\delta M^K$ in the ledger term of~\eqref{eq:sp:lagrangian}, differentiation with respect to
$I$ and to $K$ gives
\begin{align}
q &= 1-\bar\Phi^Ip^W,
&
\dot q &= \left(r+\delta\right)q-\Big[F_K-\delta\bar\Phi^Ip^W\Big],
\label{eq:app:const:capital}
\end{align}
equivalently $r=\left(F_K-\delta\bar\Phi^Ip^W\right)/q-\delta+\dot q/q$. A unit of capital now carries
its material with it, so its dividend is the marginal product net of the waste that its depreciation
generates, and its price is unity less the material liability it brings. The contrast with the baseline
is worth marking: there the materials block collapsed entirely into $q$ and the return
was $F_K\left(1-\zeta\Omega^Y\right)/q$, because the liability was carried by the separate price $p^M$;
here it is charged directly and appears in the dividend.

The two accounts agree, as they must. Writing $q^{\text{(A)}}$ for the price in
Appendix~\ref{sec:app:exogenousPhiI} with $\Phi^I$ held constant, the definitions differ by the value
of the material that a unit of capital carries,
\begin{align}
q &= q^{\text{(A)}}-\bar\Phi^Ip^M
\;=\;1-\bar\Phi^I\left(p^W-p^M\right)-\bar\Phi^Ip^M ,
\label{eq:app:const:reconciliation}
\end{align}
and substituting $\dot p^M=\left(r+\delta\right)p^M-\delta p^W$
from~\eqref{eq:app:exo:costate-M} into $\dot q=\dot q^{\text{(A)}}-\bar\Phi^I\dot p^M$
returns~\eqref{eq:app:const:capital}. The allocation is the same; what differs is whether the material
liability is priced inside $q$ or alongside it. The price $p^M$ remains available as a derived object
through~\eqref{eq:app:const:reconciliation}, but nothing requires it.

Every remaining condition is~\eqref{eq:app:exo:C}--\eqref{eq:app:exo:costate-M} with $\Phi^I=\bar\Phi^I$
and the investment condition~\eqref{eq:app:exo:I} replaced by~\eqref{eq:app:const:capital}. In particular the
recycling block, the phase structure of Proposition~\ref{prop:sp:phases} and the
threshold~\eqref{eq:sp:phases:threshold} are untouched.

\subsection{Implementation}\label{sec:app:const:market}

The policy is~\eqref{eq:app:exo:policy} with $s^I=\bar\Phi^Ip^W$, and it simplifies on the household's
side in the same way the planner's problem does. The household pays $\tau^{\circ}\delta M^K$ on the
material its capital releases and receives $s^I$ on gross investment, so with
$M^K=\bar\Phi^IK$ and $\tau^{\circ}=p^W$ its net material payment is
\begin{align}
p^W\bar\Phi^I\delta K-p^W\bar\Phi^I\left(I+\delta K\right) &= -\,p^W\bar\Phi^II :
\label{eq:app:const:hh}
\end{align}
the two instruments collapse into a single subsidy at rate $\bar\Phi^Ip^W$ per unit of \emph{net}
investment. The household no longer carries $M^K$ as a private state, and the deposit--refund of
Section~\ref{sec:mkt:implementation:reading} reduces to the statement that material locked into the
capital stock is charged only to the extent that the stock is growing.

\begin{remark}[When the reduction is legitimate]\label{rem:app:const:legit}
Constancy of $\Phi^I$ is a strong assumption in a model whose point is that material intensities move,
and it removes by construction the channel that~\eqref{eq:sp:setup:phiK-motion} exists to describe:
$\Phi^K$ can no longer drift, so the composition of the capital stock carries no information. It is
nonetheless the right variant for two purposes. It is the cheapest version to solve, hence the natural
first implementation and the natural test case for the corner-detection routine of
Section~\ref{sec:qmkt:solution}, which is where the paper's headline number lives. And it isolates the
material-storage motive in its purest form: a single parameter $\bar\Phi^I$, a single price $p^W$, and
a subsidy on net investment.
\end{remark}
